% =============================================
% REVISED BYLAWS CONTENT FILE
% =============================================
\begin{enumerate}
    \item These procedures govern the hearings held by all Institute Hearing Boards and the Review Board.
    
    \item All appeals to an Institute Hearing Board or the Review Board must be submitted through the Office of the Dean of Students. Appeal requests shall be made in writing, based on one or more of the three Institute grounds for appeal, and must specifically state the reason and rationale for the request.
    
    \item Prior to a hearing, both parties shall independently provide all facts and evidence they wish to present to the Board Chair by the specified deadline. This written evidence may include incident reports, Public Safety reports, and judicial inquiry notes. Parties must also submit a list of potential witnesses, including their full names, titles, contact information, and a justification for their relevance. The Board Chair will review these requests, approve relevant witnesses, and compile all materials into a single evidence packet. This packet will be distributed to both parties, board members, and the Senior Judicial Administrator at least 48 hours before the hearing. The report may include the student's or group's past disciplinary records for sanctioning purposes.
    
    \item A lower board may refer a case directly to the Review Board by providing a written statement of their reasons. Similarly, the Review Board may refer a case to the President.
    
    \item The Chair or their appointee shall set the date, time, and location of the hearing, notifying all board members and parties involved. Hearings will typically be held within three to ten Institute business days following the receipt of the request, though the Institute reserves the right to suspend the process during breaks or holidays. Sessions will occur on weekdays between 9:00 AM and 7:00 PM and will not exceed 2.5 hours per segment unless extreme circumstances require an exception approved by the Chair and the Senior Judicial Administrator.
    
    \item To request a postponement, a party must contact the Board Chair, Vice-Chair, or Senior Judicial Administrator at least 24 hours before the hearing. Rescheduling is granted at the Chair's discretion for extreme and unforeseeable cases. If the 24-hour notice is not met, the Board may proceed as scheduled or uphold the original decision. If the Board cannot convene, they will notify the parties and reschedule for the earliest possible time.
    
    \item All parties will receive a copy of these procedures and relevant bylaws during a meeting with the Chair. At this time, a party may submit a written petition to disqualify a student board member for good cause. Disqualification requires the member to recuse themselves or a $2/3$ vote of the Board membership.
    
    \item Hearings are closed to the public unless authorized by the Board Chair. However, victims of physical violence, including sexual offenses, are permitted to attend the entire hearing upon their request.
    
    \item The Chair rules on all questions of relevance and procedure. During the hearing, these rulings may be overruled by a $2/3$ vote of the Board if there is good cause.
    
    \item The Senior Judicial Administrator shall be present throughout the evidence review, deliberation, and decision-making process to provide guidance on procedure and respond to inquiries.
    
    \item Both parties have the right to be assisted by an officially designated Student Judicial Advisor from the Rensselaer student community. While the advisor may support the party they are assisting, they do not participate directly in the proceedings. Both parties maintain the right to testify, present the evidence packet, and cross-examine witnesses.
    
    \item At the start of the hearing, the accused party will be introduced to the board members and informed of the applicable appeal procedures.
    
    \item The hearing shall proceed in the following order:
    \begin{enumerate}
        \item The party bearing the burden of proof delivers an opening statement.
        \item The opposing party delivers their opening statement.
        \item Witnesses are called in order: first by the party bearing the burden of proof, then by the opposing party, and finally by the Board. Each witness is questioned first by the person who called them, then by the opposing party, and finally by the Board members.
        \item The Board may question either party, and the parties may question each other.
        \item The party bearing the burden of proof delivers a closing statement.
        \item The opposing party delivers their closing statement.
        \item The hearing is closed for private deliberations.
    \end{enumerate}
    
    \item The Chair will provide the Senior Judicial Administrator with the written decision within two Institute business days.
    
    \item Within three Institute business days of the hearing's conclusion, the Senior Judicial Administrator will inform the accused of the Board’s decision and recommendations, either in person or in writing.
\end{enumerate}